\begin{lemma}[Proved: directed-poset presentations]\label{cat:I7}
Every small $\kappa$-filtered category admits a final functor from a
small $\kappa$-directed poset.  Thus the poset formulation used on
161-2/060 and in Lemma~\ref{cat:I2} loses no presentations.
\end{lemma}
\begin{proof}
Let $D$ be the category.  We construct a poset $P$ with a functor
$u:P\to D$, keeping every principal downset of cardinal less than
$\kappa$.  Start with one isolated vertex for every object of $D$,
with the evident values of $u$.  At a successor stage, for each subset
$S$ of existing vertices of size less than $\kappa$, take the union of
their downsets.  This has fewer than $\kappa$ arrows, by regularity.
Choose a cocone on its image in $D$, and adjoin a new vertex above
that union with this value and these transition maps.  Do not put
relations between vertices added at the same stage.

Also, for each object $d\in D$, pair of old vertices $v,w$, and maps
$a:d\to u(v)$, $b:d\to u(w)$, adjoin an upper vertex whose transition
maps equalize these two maps.  This is possible by taking a cocone on
the union of the two downsets together with a new formal vertex of
value $d$ and arrows of values $a,b$.  The resulting diagram still
has fewer than $\kappa$ arrows.  Its cocone restricts to the old
poset diagram, so it extends $u$ coherently.  Each new principal
downset still has size less than $\kappa$.
Take unions at limit stages and repeat for $\kappa$ stages.
All stages are sets.  Every subset of fewer than $\kappa$ vertices
appears before a common stage and gets an upper bound at the next
one, so $P$ is $\kappa$-directed.  For each $d$, the comma category
$d/u$ is nonempty by the initial vertices, and any two of its objects
are joined by the added equalizing upper vertex.  Hence it is connected.
A cocone on $u$ therefore gives a well-defined map from each $d$,
independent of the choice of an object in $d/u$; the equalizing
vertices verify both independence and compatibility with every arrow
of $D$.  This is exactly finality, directly from the colimit universal
property.
\end{proof}

\begin{theorem}[Proved: the cofinal-rank replacement for the false corollary]\label{cat:I8}
Let $E,F$ be $\kappa$-accessible and let $T:E\to F$ preserve
$\kappa$-filtered colimits.  There are arbitrarily large regular
cardinals $\lambda$ such that $E,F$ are $\lambda$-accessible and
$T$ preserves $\lambda$-filtered colimits and
$\lambda$-presentable objects.
More explicitly, choose small full categories of $\kappa$-presentable
generators $C,D$ for $E,F$, respectively, and for every $c\in C$
choose a $\kappa$-directed presentation of $T(c)$ by objects of $D$.
If $\nu\geq\kappa$ bounds all those indexing sets, any
$\mu\geq\nu$ with $\mu^{<\kappa}=\mu$ and
$\lambda=\mu^+$ has the asserted property.
In particular one may take $\mu=2^\nu$ after increasing $\nu$.
\end{theorem}
\begin{proof}
By Lemma~\ref{cat:I7}, write any object of $E$ as a colimit of a
$\kappa$-directed poset of objects of $C$.
Lemmas~\ref{cat:I2} and~\ref{cat:I3} express it as a
$\lambda$-directed colimit of the colimits of its
$\kappa$-directed subdiagrams of size at most $\mu$.
Each such small colimit is $\lambda$-presentable: its Hom functor
is a limit of at most $\mu<\lambda$ Hom functors of
$\kappa$-presentables, hence of $\lambda$-presentables, and those
limits commute with $\lambda$-filtered colimits of sets.
There is only a set of isomorphism types of such small diagrams.
This proves $\lambda$-accessibility of $E$, and the same proof
applies to $F$.

If $X$ is $\lambda$-presentable, its identity factors through one
stage in this canonical $\lambda$-directed presentation.  Thus it is
a retract of a colimit of a $\kappa$-directed diagram of at most
$\mu$ objects of $C$.  By the bound on the chosen presentations,
every $T(c)$ is $\lambda$-presentable, using the same Hom-limit
argument in $F$.  Applying $T$ to that small directed diagram gives
its colimit, which is again $\lambda$-presentable.  Since a retract
of a presentable object is presentable, $T(X)$ has the required rank.
A $\lambda$-filtered category is $\kappa$-filtered, so preservation
of its colimit already follows from the hypothesis on $T$.
Finally $(2^\nu)^\alpha=2^\nu$ for $0<\alpha<\kappa\leq\nu$,
which verifies the indicated bound and gives arbitrarily large choices.
This replaces ``every sufficiently large rank'' by a proved cofinal
family of ranks.
\end{proof}

\begin{proposition}[Proved: the exact uniform cardinal condition]\label{cat:I9}
For infinite regular cardinals $\kappa<\lambda$, the following are
equivalent:
\begin{enumerate}
\item For every set $X$ of size less than $\lambda$, the poset
$[X]^{<\kappa}$ has a cofinal subset of size less than $\lambda$.
\item In every $\kappa$-directed poset, every subset of size less than
$\lambda$ is contained in a $\kappa$-directed subset of size less than
$\lambda$.
\end{enumerate}
This condition is often written $\kappa\mathrel{\triangleleft}\lambda$.
It is the exact uniform closure condition; a power equality is a
convenient sufficient cardinal bound, not its definition.
The unrestricted closure step used on scan 061 is false for
$\kappa=\aleph_1$ and allowed subdiagram size $\mu=\aleph_\omega$.
\end{proposition}
\begin{proof}
Assume (1) and begin with a subset $A$ of a directed poset of size
less than $\lambda$.  Choose a cofinal family in $[A]^{<\kappa}$
with fewer than $\lambda$ members and adjoin upper bounds only for
members of that family.  Repeat for $\kappa$ stages.  Regularity of
$\lambda>\kappa$ keeps the cardinality less than $\lambda$.
Regularity of $\kappa$ puts every fewer-than-$\kappa$ family of
resulting elements in one stage; the chosen cofinal family at that
stage supplies its upper bound.  This proves (2).
Conversely, apply (2) to the $\kappa$-directed poset
$[X]^{<\kappa}$ and its subset of singletons.
Any $\kappa$-directed subposet containing those singletons is
cofinal, because every member of $[X]^{<\kappa}$ is the union of
fewer than $\kappa$ singletons.  This proves (1).

For the counterexample let $I=[\mu]^{\leq\aleph_0}$, ordered by
inclusion.  It is $\aleph_1$-directed.  If a subposet $J$ of size
at most $\mu$ contains all singletons and is $\aleph_1$-directed,
then it is cofinal in $I$ by the same union argument.
But no family $(A_i)_{i<\mu}$ of countable subsets is cofinal.
Partition its indices into countably many sets $I_n$, each of size
less than $\mu$; this is possible because $\operatorname{cf}\mu=\omega$.
Choose $a_n\in\mu\setminus\bigcup_{i\in I_n}A_i$, since that union
has cardinal less than $\mu$.  The countable set $\{a_n:n\in\mathbf N\}$
is contained in none of the $A_i$, a contradiction.
Thus the missing cardinal hypothesis is real.  Lemma~\ref{cat:I2}
provides the corrected closure estimate, while (1) gives its exact
uniform form.
\end{proof}

